-----Step 9: cumulative residual sd plot-----

1. Purpose

The cumulative residual sd plot answers the question:

   What is a good fitted model for the data?

The prior 8 steps in this analysis sequence
addressed the important 2 goals regarding:

   1. Factors: determining the most important factors
      which affect the response, and

   2. Settings: determining the best settings for these
      factors.

In addition to the above, a third goal is of
interest:

   3. Model: determining a fitted model (that is, a
      prediction equation) which functionally
      relates the observed response Y with the
      various main factors and interactions.

Such a function makes particular sense when all of the
individual factors are continuous and ordinal (such as
temperature, pressure, humidity, concentration, etc.)
as opposed to any of the factors being discrete and
non-ordinal (such as plant, operator, catalyst,
supplier).

In the former continuous-factor case, the analyst would
use such a function for

   1. Reproduction/Smoothing: predict the response
      at the observed design points;

   2. Interpolation: predict what the response would be
      at (unobserved) regions between the design
      points; or

   3. Extrapolation: predict what the response would be
      at (unobserved) regions beyond the design
      points.

For the discrete-factor case, the methods
developed below to arrive at such a function still
apply, and so the resulting fitted model may be
used for reproduction, but obviously the
interpolation and extrapolation aspects do not
apply.

In modeling, therefore, we thus seek a function f in
the k factors X1, X2, ..., Xk such that the
post-fitting predicted values "Yhat"

   Yhat = f(X1,X2,...,Xk)

are "close" to the observed raw data values Y. To
this end, 2 tasks exist:

   1. Determine a good functional form f;

   2. Determine good estimates for the coefficients in
      that function f.

Thus, for example, if we had 2 factors X1 and X2,
then our goal would be to

   1. determine some function Y = f(X1,X2); and

   2. estimate the parameter coefficients in f

such that the resulting fitted model would
yield predicted values (Yhat) that are as close as
possible to the observed response values Y.
If the form f has been properly chosen, and the model
coefficients have been estimated well, then a good
model will result and that model will have the
characteristic that the differences ("residuals" =
Y-Yhat) will be uniformly near-zero.  On the other
hand, a poor model (from a poor choice of the form f,
or a poor job of estimating the model parameters) will
have the characteristic that (some or all) of the
residuals Y-Yhat will be "large".

For a given model, a single statistic which summarizes
the quality of a fit via the typical size of the n
residuals is the residual standard deviation:

      Sres = sqrt(...)

If all such residuals are small (that is, we have
a good-fitting model), then Sres will be small; if
some or all of the residuals are large (that is,
we have a poor-fitting model), then Sres will be
large.

For a given data set, each proposed model has its own
quality of fit, and hence its own residual standard
deviation.  Clearly, the residual standard deviation is
more of a model-descripter than a data-descripter.
Whereas "nature" creates the data, the analyst
creates the models--theoretically, for the same
data set, it is possible for the analyst to
propose an indefinitely large number of models.

In practice, however, an analyst will usually
foward only a small, finite number of plausible
models for consideration.  Each such model will
have its own residual standard deviation.  The
cumulative residual standard deviation plot is
simply a graphical representation of this
collection of residual standard deviations for
various models.  The plot is beneficial in that

   1. good models are distinguished from bad models;

   2. simple good models are distinguished from
      complicated good models.

In summary, then, the cumulative residual sd plot is a
graphical tool to help assess

   1. which models are poor (least desirable); and

   2. which models are good but complex (more desirable); and

   3. which models are good and simple (most desirable).

2. Output

   1. Primary  : A good-fitting prediction
                 equation consisting of an
                 additive constant plus the most
                 important main effects and
                 interactions.

   2. Secondary: The residual standard deviation
                 for this good-fitting model.

3. Definition

   1. Vertical Axis    : Ordered (largest to
                         smallest) residual
                         standard deviations of a
                         sequence of progressively
                         more complicated fitted
                         models.


   2. Horizontal Axis  : Factor/interaction
                         identification of the last
                         term included into the cumulative
                         linear combinatoric model:
                             1 indicates factor X1;
                             2 indicates factor X2;
                             ...
                             12 indicates the 2-term
                                X1*X2 interaction
                             123 indicates the 3-term
                                X1*X2*X3 interaction
                             etc.

   3. Far right margin : Factor/interaction
                         identification (built-in
                         redundancy):

                             1 indicates factor X1;
                             2 indicates factor X2;
                             ...
                             12 indicates the 2-term
                                X1*X2 interaction
                             123 indicates the 3-term
                                X1*X2*X3 interaction
                             etc.

                         If the design is a
                         fractional factorial,
                         then the confounding
                         structure is provided for
                         main factors and 2-term
                         interactions.

The cumulative residual standard deviations plot
is thus a Pareto-style--largest to smallest--
graphical summary of residual standard deviations
for a selected series of progressively more
complicated linear combinatoric models.

A linear combinatoric model is a function which starts
with an additive constant and then sequentially
augments the model by including--one at a time--
remaining factors and interactions.  Each factor and
interaction is incorporated into the model in an
additive (rather than in a multiplicative or
logarithmic or power, etc. fashion).  At any stage, the
ordering of the next term to be added to the model is
such that it will result in the maximal decrease in the
resulting residual standard deviation.

4. Motivation:

This section addresses the following 9 questions:

   1. What is a model?
   2. How do we construct a goodness-of-fit metric for a model?
   3. How do we construct a good model?
   4. How do we know when to stop adding terms?
   5. What is the final form for the model?
   6. Why is the 1/2 in the model?
   7. What are the advantages of the linear combinatoric model?
   8. How do we use the model to generate predicted values?
   9. How do we use the model beyond the data domain?
  10. What is the best confirmation point for interpolation?
  11. How do use the model for interpolation?
  12. How do use the model for extrapolation?

..............................................
Question 1: What is a model?

A model is a mathematical functions which relates
the response Y to the factors X1 to Xk.  A model
has a

   1. functional form; and
   2. coefficients.

An excellent and easy-to-use functional form that
we find particularly useful is a linear
combination of the main factors and the
interactions (of all orders).  As for the
coefficients in this linear combinatoric model,
high-quality, easy-to-obtain coefficient estimates
are gotten via application of the least squares
estimation criterion (regression).  A given
functional form with estimated coefficients is
referred to as a "fitted model" or a "prediction
equation"

For given settings of the factors X1 to Xk, a
fitted model will yield predicted values.  For
each (and every) setting of the Xi's, a
"perfect-fit" model is one in which the yielded
predicted values are identical to the observed raw
response Y's at these Xi's.  In other words, a
perfect-fit model will yield a vector of predicted
values identical to the observed vector of
response values.  For these same Xi's (each and
all), a "good-fitting" model is one which yields
predicted values "acceptably near"--but not
necessarily identical to--the observed responses Y.

The residuals (= deviations = error) of a model
are the vector of differences (Y - Yhat) between
the raw responses and the predicted values from
the model.  For a perfect-fit model, the vector of
residuals will be all zeros.  For a good-fitting
model, the vector of residuals will be acceptably
(from an engineering point of view) close to zero.


..............................................
Question 2: How do we construct a goodness-of-fit metric
            for a model?

This question deals with the issue of how to
construct a metric--a statistic--which may be used
to ascertain the quality of the fitted model.  The
statistic should be such that for one range of
values, the implication is that the model is good,
whereas for another range of values, the
implication is that the model is poor-fitting.

Since a model's adequacy is inversely related to
the size of its residuals, then one such obvious
statistic is

   the sum of the absolute residuals.

Clearly, the smaller this sum, the smaller the
residuals, which implies the closer the predicted
values (to the raw data Y), and hence the better
the fitted model.  A disadvantage of this
statistic is that it may grow larger, simply as
the sample size n grows larger.

A better metric which does not change (much) with
increasing sample size is the average absolute
residual:

   AAR = (sum of absolute residuals) / n

where n is the number of response values.  Again,
small values for this statistic imply
better-fitting models.

An alternative, but similar, metric which has
better statistical properties is

   sqrt((sum of squared residuals) / n)

As with the previous, the smaller this statistic,
the better the model.

Our final metric which is to be preferred because
of its excellent statistical properties
(unbiasedness) is

   sqrt((sum or squared residuals) / (n-p))

where p is the number of fitted coeffients in the
model.  This final statistic is called the
"residual standard deviaion--it is the standard
deviation of the residuals from a given model.
The smaller this residual standard deviation, the
better-fitting the model.  We shall use the
residual standard deviation as our metric of
choice for evaluating and comparing various
proposed models.


..............................................
Question 3: How do we construct a good model?

Given that we have a statistic to measure the
quality of a model--any model--we move to the
question as to how to construct reasonable models
for fitting data resulting from 2**k and 2**(k-p)
designs.

<Initial simple model>
The simplest such proposed model is

   Y = f(X1,X2,...,Xk) = c + e

that is, the response Y = a constant + random
error.  This trivial model says that all of the
factors (and interactions) are in fact worthless
for prediction and so the best-fit model is one
which consists of a simple horizontal straight
line through the body of the data.  The least
squares estimate for this constant c in the above
model is the sample mean Ybar:

   estimated c = Ybar = average Y

The prediction equation for this model is thus

   Yhat = Ybar

The predicted values Yhat for this fitted trivial
model is thus a vector consisting of the same
value (namely Ybar) throughout.  The residual
vector for this model will thus consist of simple
deviations about the mean:

   Y - Ybar

Since the number of fitted coefficients in this
model is 1 (namely the constant c), then the
residual standard deviation becomes the following:

   sqrt((sum of squared residuals about Ybar) / (n-1))

which is of course identically the familiar and
commonly-employed sample standard deviation.  If
the residual standard deviation for this trivial
model is "small enough", then we may terminate the
model-building process right there with no further
inclusion of terms.  In practice, however, this
trivial model does NOT usually yield a residual
standard deivation which is small enough (because
the common value Ybar is typically not acceptably
close enough to some of the raw responses Y) and
so the model must be simply augmented--but how?

<Next-step model>
The logical next-step proposed model will consist
of the above additive constant + plus some term
which will improve the predicted values the most
(and equivalently will reduce the residuals the
most, and thus reduce the residual standard
deviation the most).

<Using the most important effects>
As it turns out, it is a mathmatical fact that the
the factor (or interaction) which has the largest
estimate effect E (= Ybar(+) - Ybar(-)) will
necessarily-- after being included in the
model--yield the "biggest bang for the buck" in
terms of improving the predicted values toward the
response Y. Hence at this point the model-buiding
process and the effect estimation process merge.

From all of the previous 8 steps in our analysis,
and based on the magnitude of the estimated
effects, we developed a ranked list of factors and
interactions-- most important to least important.
We thus have a ready-made ordering of the terms
which could be added--one at a time--to the model.
This ranked list of effects is precisely the
"prescription" we need to cumulatively build more
complicated--but better fitting--models.

Our procedure will thus be to step down--one by
one--the ranked list of effects, cumulatively
augmenting our current model by the next term in
the list, and then compute (for all n design
points) predicted values, residuals, and residual
standard deviation.  We continue this
one-term-at-a-time augmentation until the
predicted values are acceptably close to the
observed responses Y (and hence the residuals and
residual standard deviation become acceptably
close to zero).

Starting with the simple average, each cumulative
model in this iteration process will have its own
associated residual standard deviation.  In
practice, the iteration continues until the
residual standard deviations become sufficiently
small.

The cumulative residual standard deviation plot is
a graphical summary of the above model-building
process.  On the horizontal axis is a series of
terms (starting with the average, and continuing
on with various main factors and interactions).
After the average, the ordering of terms on the
horizontal axis is identical to the ordering of
terms based on step 8's ranking based on effect
magnitude.

On the vertical axis is the corresponding residual
standard deviation that results when the
cumulative model has its coefficients fitted (via
least squares), and then has its predicted values,
residuals, and residual standard deviations
computed.  The first residual standard deviation
(on the far left of the cumulative residual
standard deviation plot) is that which results
from the model consisting of

   1) the average.

The second residual standard deviation plotted is
from the model consisting of

   1) the average, plus
   2) the term with the largest |effect|.

The third residual standard deviation plotted is
from the model consisting of

   1) the average, plus
   2) the term with the largest |effect|, plus
   3) the term with the second largest |effect|.

and so forth.

..............................................
Question 4: How do we know when to stop adding terms?

Proceeding left to right, as we add more terms to
the model, the cumulative residual standard
deviation "curve" will typically decrease.  At the
beginning (on the left), as we add large-effect
terms, the decrease from one residual standard
deviation to the next residual standard deviation
will be large.  The incremental improvement
(decrease) then tends to drop off slightly.  At
some point the incremental improvement will
typically slacken off considerably.
Appearance-wise, it is thus very typical for such
a curve to have a "hockey stick" appearance:

   1. starting with a series of large decrements
      between successive residual standard
      deviations; then

   2. hitting an elbow; then

   3. having a series of gradual incremental
      decreases thereafter.

<Stopping rule>
The cumulative residual standard deviation plot
provides a visual answer to the question:

   What is a good model?

by answering the related question:

   When do we stop adding terms to the cumulative model?

Interpretationally, the most common stopping rule
for adding terms is to cease immediately upon
encountering the "elbow".  We include all terms up
to (and including) the elbow point (after all,
each of these terms decreased the residual
standard deviation by a large amount); but we
exclude any terms afterward (since every
successive term thereafter does not decrease the
residual standard deviation fast enough to warrant
inclusion in the model).

..............................................
Question 5: What is the final form for the model?

From the above discussion, we thus note and
recommend a final form for the model which
consists of an additive constant plus a linear
combination of main effects and interactions.
What then is the specific form for the linear
combination?  We recommend the following:

   For the trivial k = 1 factor case:

      Y = f(X1) = c + (1/2)*(B1*X1)

   For the k = 2 factor case:

      Y = f(X1,X2) = c + (1/2)*(B1*X1 + B2*X2 + B12*X1*X2)

   For the k = 3 factor case:

      Y = f(X1,X2,X3) = c + (1/2)*(B1*X1 + B2*X2 + B3*X3 +
                                   B12*X1*X2 + B13*X1*X3 +
                                   B23*X2*X3 + B123*X1*X2*X3)

and for the general k case:

      Y = f(X1,X2,...,Xk) = c + (1/2)*(linear combination
                                       of all main factors
                                       and all interactions
                                       of all orders)

<Ordered linear combination>
Although the listing above has the terms ordered
with the main factors, and then the 2-term
interactions, then the 3-term interactions, etc., in
practice, it is recommended that the terms be
re-ordered so that the most important terms
(whether they be main effects or interactions) be
listed first, and then other terms listed in
decreasing order of importance thereafter).  Aside
from providing a functional representation of the
response, models should help to functionally
reinforce what is driving the response, and so
such a re-ordering helps to assist in that regard.
Thus for k = 2, if factor 2 is most important, and
the 2-term interaction is next in importance, and
factor 1 is least important, then it is
recommended that the above ordering of

   Y = f(X1,X2) = c + (1/2)*(B1*X1 + B2*X2 + B12*X1*X2)

be rewritten as

   Y = f(X1,X2) = c + (1/2)*(B2*X2 + B12*X1*X2 + B1*X1)


..............................................
Question 6: Why is the 1/2 in the model?

The leading 1/2 is a multiplicative constant that
we have chosen to include in our expression of the
linear combinatoric model.  The presence or
absence of the arbitrary 1/2 term does NOT affect
the predictive quality of the model after least
squares fitting.  Clearly, if we choose to exclude
the 1/2, then the least squares fitting process
will simply yield estimated values of the
coefficients B that are twice the size of the
coefficients B that would result if we included
the 1/2.

We recommend the inclusion of the 1/2 because of
an additional property that we would like to
impose on the model, namely, we desire that:

   The value of the least squares estimated
   coeffient B for a given factor (and
   interaction) be visually identical to the
   estimated effect E for that factor (and
   interaction).

For a given factor X2 (say), the estimated least
squares coefficient B2 and the estimated effect
(E2) are NOT in general identical--either in value
or in concept.

<Effect>
For factor X2, the "effect" E2 is defined as the
change in the mean response as we proceed from the
- setting of the factor to the + setting of the
factor, symbolically:

   estimated effect = Ybar(+) - Ybar(-)

Note that the estimated effect E2 value does NOT
involve a model per se, and is definitionally
invariant to any other factors and interactions
that may effect the response.  We examined and
derived the factor effects E in steps 1 to 8 of
the general DEX analysis procedure.

On the other hand, the estimated coefficient B2 in
a model is defined as that value which results
when we place the model into the least squares
fitting process (regression).  The value that
returns for B2 depends--in general--on the form
for the model, on what other terms are included in
the model, and on the experiment design that was
run.  The least squares estimate for B2 is mildly
complicated--involving a behind-the-scenes design
matrix multiplication and inversion.  The
coefficient values B that result are generally
obscured by the mathematics to make the
coefficients have the collective property that the
fitted model as a whole yield a minimum sum of
squared deviations ("least squares").

<Orthogonality>
Rather remarkably, these 2 concepts and values:

   1. factor and interaction effect estimates E, and

   2. least squares coefficient estimates B

merge for the class of experiment designs for
which this 10-step procedure was developed,
namely, 2-level full and fractional designs which
are orthogonal.  Orthogonality has been promoted
and chosen because of its desirable design-side
properties (every factor is balanced and every
2-term cross-product is balanced).  But to boot,
orthogonality has 2 extraordinary properties on
the data analysis side:

   1. For the above linear combinatoric
      models, the usually complicated least
      squares estimates for the coefficients B
      (which generally involves matrix
      multiplication and inversion) reduces to a
      computationally trivial and familiar form,
      namely,

            B = E = Ybar(+) - Ybar(-)

   2. The usual general modeling property that the
      least squares estimate for a factor
      coefficient changes depending on what other
      factors have been included in (or excluded
      from) the model is now moot.  With
      orthogonal designs, the coefficient
      estimates are invariant-- the estimate
      (e.g., B2) for a given factor (e.g., X2)
      will NOT change as other factors and
      interactions are included in (or excluded
      from) the model; that is, the estimate for
      the factor 2 effect (B2) remains the same
      regardless of what other factors are
      included in the model.

The net effect of the above two properties is that
a factor effect can be computed once and for all,
and that value will hold for any linear
combinatoric model involving that term--
regardless of how simple or how complicated the
model.  This process greatly simplifies the
model-building process because the need is thus
eliminated to recalculate all of the model
coefficients for each new model that gets created
as new terms are added to or dropped.

Given the calculational simplicity of orthogonal
designs, why then is 1/2 the appropriate
multiplicative constant?  Why not 1/3, 1/4, etc.
To answer this, we revisit our specified desire
that

   when we view the final fitted model
   and look at the coefficient associated with X2,
   say, then it is desirable that the value of that
   coefficient B2 reflect identically the expected
   total change delta Y in the response Y as we
   proceed from the - setting of X2 to the +
   setting of X2 (that is, we would like the
   estimated coefficient B2 to be identical to the
   estimated effect E2 for factor X2).

Thus in glancing at the final model in this form,
the coefficients B of the final model will
immediately reflect not only the relative
importance of the coefficients one to another, but
will also reflect (absolutely) the effect of the
associated term (main effect or interaction) on
the response.

In general, the least squares estimate of a
coefficient in a linear combinatoric model will
yield a coefficient which is essentially a slope:

   delta Y / delta X

associated with a given factor X. Thus in order to
achieve the desired interpretation of the
coefficients B as being the raw change in the Y
(delta Y), then we must account for and remove the
change in X (delta X).

What is the delta X? In our design descriptions,
we have chosen the notation of Box, Hunter and
Hunter (1978) and set each (coded) factor to
levels of - and +. This - and + is a shorthand
notation for -1 and +1.  The advantage of this
notation is that 2-term interactions (and any-term
interaction) also uniformly takes on the closed
values of -1 and +1, since

   -1*-1 = +1
   -1*+1 = -1
   +1*-1 = -1
   +1*+1 = +1

and hence the set of values that the 2-term
interactions (and all interactions) take on are
{-1,+1}.  This -1 and +1 notation is superior in
its consistency than the 1,2 notation of, say,
Taguchi in which case the interaction X1*X2, say,
would take on the values

   1*1 = 1
   1*2 = 2
   2*1 = 2
   2*2 = 4

which yields the set {1,2,4}.  Hence, with the
-1,+1 values for the main factors, we also have
-1,+1 values for all interactions which in turn
yields (for all terms) a consistent delta X of

   delta X = (+1) - (-1) = +2

In summary then,

   B = (delta Y) / (delta X) = (delta Y) / 2
                             = (1/2) * (delta Y)

and so to achieve our goal of having the final
coefficients reflect delta Y only, we simply
gather up all of the 2's in the denominator and
create a leading multiplicative constant of 1 with
denominator 2, that is, 1/2.

For example, for the trivial k = 1 case, the
obvious model

   Y = intercept + slope*X1

   Y = c + ((delta Y)/(delta X))*X1

 becomes

   Y = c + (1/delta X) * (delta Y)*X1

or simply

   Y = c + (1/2) * (delta Y)*X1

   Y = c + (1/2)*(factor 1 effect)*X1

This k = 1 factor result is easily seen to extend
to the general k case.

..............................................
Question 7: What are the advantages of the linear
            combinatoric model?

The recommended linear
combinatoric model consisting of main factors and all
interactions has 2 advantages:

   1. Perfect Fit:  If we choose to include in the
                    model all of the main factors
                    and all interactions (of all
                    orders), then the resulting
                    least squares fitted model
                    will have the property that
                    the predicted values will be
                    IDENTICAL to the raw response
                    values Y. (We will illustrate
                    this in the next section.)

   2. Comparable Coefficients:  Since the model fit has
                    been carried out in the coded
                    factor (-1,+1) units rather
                    than the units of the original
                    factor (temperature, time,
                    pressure, catalyst
                    concentration, etc.), then the
                    factor coefficients
                    immediately become comparable
                    to one another, which serves as
                    an immediate mechanism for the
                    ranking of the relative
                    importance of the factors--
                    scale-free.

To illustrate in detail the above latter point,
suppose the (-1,+1) factor X1 is really a coding
of temperature T (with the orginal temperature
ranging from 300 to 350 degrees), and suppose the
(-1,+1) factor X2 is really a coding of time t
(with the original time ranging from 20 to 30
minutes).  Given that, then a linear combinatoric
model in the original temperature T and time t
would yield coefficients whose magnitude depends
on the magnitude of T (300 to 350) and t (20 to
30), and whose value would change if we decided to
change the units of T (from fahrenheit degrees to
centigrade degrees) and t (from minutes to
seconds).  All of this is avoided by carrying out
the fit not in the original units for T (300,350)
and t (20,30), but in the coded units of X1
(-1,+1) and X2 (-1,+1).  The resulting
coefficients become unit-invariant, and thus the
coefficient magnitudes reflect the true
contribution of the factors and interactions,
without regard to unit.

Such coding leads to no loss of generality, since
the coded factor may be expressed as a simple
linear relation to the original factor (X1 to T,
X2 to t), and so the unit-invariant coded
coefficients may be easily transformed to
unit-sensitive original coefficients if so
desired.

..............................................
Question 8: How do we use the model to generate
            predicted values?

To illustrate the details as to how a model may be
used for prediction, let us consider a simple case
and generalize from it.  Consider the simple
Yates-order k = 2 factor 2**2 full factorial
design in X1 and X2:

        X1     X2
         -      -
         +      -
         -      +
         +      +

We augment this design with a simple (for
illustration purposes) response vector Y:


        X1     X2      Y
         -      -      2
         +      -      4
         -      +      6
         +      +      8

or geometrically:

       +1 --------------
          |6          8|
          |            |
          |            |
     X2   |            |
          |            |
          |            |
          |2          4|
       -1 --------------
          -1          +1
                X1

For this k = 2 case, the recommended model as
discussed before would be

   Y = c + (1/2)(B1*X1 + B2*X2 + B12*X1*X2)

From the above diagram, we may deduce that the
estimated factor effects must be:

   c   = the average response = Ybar
       = (2 + 4 + 6 + 8) / 4 = 5

   B1  = average change in the response Y as X1
         goes from -1 to +1
       = ((4-2) + (8-6)) / 2 = (2 + 2) / 2 = 2

   B2  = average change in the response Y as X2
         goes from -1 to +1
       = ((6-2) + (8-4)) / 2 = (4 + 4) / 2 = 4

   B12 = interaction = (the less obvious) average change
         in the response Y as X1*X2 goes from -1 to +1
       = ((2-4) + (8-6)) / 2 = (-2 + 2) / 2 = 0

and so the fitted model (that is, the prediction
equation) is

   Yhat = 5 + (1/2)*(2*X1 + 4*X2 + 0*X1*X2)

or preferably (with the terms rearranged in
decending order of importance)

   Yhat = 5 + (1/2)(4*X2 + 2*X1 + 0*X1*X2)

Continuing with the data set and model in hand,
what predicted values Yhat does this model yield
at the 4 design points:

        X1     X2      Y      Yhat
         -      -      2       ?
         +      -      4       ?
         -      +      6       ?
         +      +      8       ?

For the first design point (X1 = -, X2 = -), we
note that the raw response is Y = +2, but what is
the predicted value?  To compute this value, note
that - and + are shorthand for -1 and +1, and so
the desired predicted value Yhat is

   Yhat(X1 = -1, X2 = -1) =
   Yhat = 5 + (1/2)(4*X2 + 2*X1 + 0*X1*X2)
   Yhat = 5 + (1/2)(4*(-1) + 2*(-1) + 0*(-1)(-1))
   Yhat = 5 + (1/2)(-6)
   Yhat = 5 - 3
   Yhat = 2

Thus for this case and this model, the predicted
value (Yhat = 2) is identical to the response
value (Y = 2):

        X1     X2      Y      Yhat
         -      -      2       2
         +      -      4
         -      +      6
         +      +      8

Similarly, for the second design point (X1 = +,
X2= -), the predicted value from this model is

   Yhat(X1 = +1, X2 = -1) =
   Yhat = 5 + (1/2)(4*X2 + 2*X1 + 0*X1*X2)
   Yhat = 5 + (1/2)(4*(-1) + 2*(+1) + 0*(+1)(-1))
   Yhat = 5 + (1/2)(-2)
   Yhat = 5 - 1
   Yhat = 4

and for the third design point:

   Yhat(X1 = -1, X2 = +1) =
   Yhat = 5 + (1/2)(4*X2 + 2*X1 + 0*X1*X2)
   Yhat = 5 + (1/2)(4*(+1) + 2*(-1) + 0*(-1)(+1))
   Yhat = 6

and for the last design point:

   Yhat(X1 = +1, X2 = +1) =
   Yhat = 5 + (1/2)(4*X2 + 2*X1 + 0*X1*X2)
   Yhat = 5 + (1/2)(4*(+1) + 2*(+1) + 0*(+1)(+1))
   Yhat = 8

Thus the k = 2 models consisting of all (= 2) main
factors and all (= 1) interactions has the property
that the predicted values Yhat are identical to
the raw response values Y:

        X1     X2      Y     Yhat
         -      -      2       2
         +      -      4       4
         -      +      6       6
         +      +      8       8

<Perfect fit>
This is a perfect-fit model.  Such perfect fit
models will result anytime (in this orthogonal
2-level design family) when we include all main
effects and all interactions of all orders.
Remarkably, this is true not only for k = 2
factors, but is is true for general k.

<Residuals>
For a given model (any model), the difference
between the original response value Y and the
predicted value Yhat is referred to as the
"residual":

   residual = Y - Yhat

The perfect-fit full-blown (all main factors and
all interactions of all orders) combinatoric
models will have all residuals identically zero.

The perfect fit is a mathematical property that comes
(free!) if we choose to use the recommended linear
combinatoric model.  Thus linear combinatoric
models have the distinct advantage of yielding a
perfect fit at the data points.

<Price for perfect fit>
What price is paid for thjis perfection?  The
obvious price is that to achieve this, we must
compute and carry the (average and the)
coefficients of all main factors and all
interactions of all orders.  Including the
average, there will in general be 2**k
coefficients to fully describe the fitting of the
n = 2**k points.  This is very much akin to the Y
= f(X) polynomial fitting of n distinct points--
it is well known that this may be done "perfectly"
by fitting a polynomial of degree n-1.  It is
comforting to know that such perfection is
mathematically attainable, but in practice do we
want to do such all the time (or anytime)?  The
answer is frequently "no"--for 2 reasons:

   1. Noise: It is very common that the response
             data Y has noise (= error) in it.  Do
             we want to go out of out way to fit
             such noise?  Or do we want our model
             to filter out the noise and just fit
             the "signal"?  For the latter, fewer
             coefficients may be in order, in the
             same spirit that we may forego a
             perfect- fitting (but jagged) 11-th
             degree polynomial to 12 data points,
             and opt out instead for an imperfect
             (but smoother) 3rd degree polynomial
             through the same 12 points.

   2. Parsimony: For full factorial designs, to fit
             the n = 2**k points, we would need to
             compute 2**k coefficients.  We gain
             information by noting the magnitude
             and sign of such coefficients, but
             numerically we have n data values Y
             as input and n coefficients B as
             output, and so no numerical reduction
             has been achieved--we have simply
             traded one set of n numbers (the
             data) for another set of n numbers
             (the coefficients).  Not all of these
             coefficients will be equally
             important-- some will be much more
             important than others.  At times that
             importance gets clouded by the sheer
             "volume" of the n = 2**k
             coefficients.  Parsimony suggests
             that our result should be simpler and
             more focused than our n starting
             points.  Hence fewer retained
             coefficients are called for.

The net result therefore is that in practice, we
may voluntarily give up the perfect--but
unwieldy--model for the imperfect--but
parsimonious--model.

<Imperfect fit>
The above calculations illustrated the computation
of predicted values for the complete model (all
main factors and all interactions).  On the other
hand, as discussed above, it may be convenient for
signal- or parsimony- purposes to deliberately
omit some (unimportant) factors.  If the analyst
chooses such a model, we note that the methodology
for computing predicted values Yhat is precisely
the same.  In such case, however, the resulting
predicted values Yhat will in general NOT be
identical to the original response values Y--that
is, we no longer get a perfect fit.  Thus, linear
combinatoric models which omit some terms will
usually have some residuals non-zero.

..............................................
Question 9: How do we use the model beyond the
            data domain?

The previous section illustrated how to compute
predicted values at the points included in the
design.  One of the virtues of modelling is that
the resulting mathematical equation (prediction
equation) is NOT restricted to the design data
points.  From the prediction equation, predicted
values can be computed elsewhere and anywhere:

   1. within the domain of the data (interpolation);

   2. outside of the domain of the data (extrapolation).

In the hands of an expert scientist/engineer/
analyst, the ability the predict elsewhere is
extremely valuable.  Based on the fitted model, we
thus have the ability to compute predicted values
for the response at an infinite number of
internal and external points, and thus the analyst
can go beyond the handful of response values Y
that he/she has at hand, but rather can get a feel
(typically via subsequent contour plotting) as to
what the nature of the entire response surface is.

This added insight into the nature of the response
is "free" and is an incredibly important benefit
of the entire model-building exercise.

<Caution>
Can we be fooled and misled by such a mathematical
and computational exercise?  After all, is not the
only thing that is "real" is the data and
everything else is artifactual?  The answer is
"yes" and so such interpolation/extrapolation is a
2-edged sword that must be wielded with care.  The
best attitude--and especially for extrapolation--
is that the derived conclusions must be viewed
with extra caution.

By construction, the recommended fitted models
will be good at the design points.  If the
full-blown model is used, the fit will be perfect.
If the full-blown model is reduced just a bit,
then the fit will still typically be quite good.
By continuity one would expect perfection/goodness
at the design points would lead to goodness in the
immediate vicinity of the design points.  However,
such local goodness does NOT guarantee that the
derived model will be good at some distance from
the design points.

<Confirmation runs>
Modeling and prediction allows us to go "beyond"
the data to gain additional insights, but such
must be done with great caution.  Interpolation is
generally safer than extrapolation, but
mis-prediction, error, and misinterpretation is
liable to occur occur in either case.

The analyst should definitely do the
model-building process and enjoy the ability to
predict elsewhere, but the analyst must always be
prepared to validate the interpolated and
extrapolated predictions by collection of
additional real, confirmatory data.  The general
empirical model that we recommend knows "nothing"
about the engineering, physics, or chemistry
surrounding your particular measurement problem,
and although the model is the best generic model
available, it must nonetheless be confirmed by
additional data.  Such additional data can be done
pre-emptively or post-experimentally.  If done
pre-emptively, a recommended procedure for
checking the validity of the fitted model is to
augment the usual 2**k or 2**(k-p) designs with
additional interpolation points--at tne center of
the design (to be discussed later).

<Continuous factors>
Of course all such discussion of interpolation and
extrapolation make sense only in the context of
continuous ordinal factors (such as temperature,
time, pressure, size, etd.) Interpolation and
extrapolation makes no sense for discrete
non-ordinal facors (such as supplier, operators,
design types, etc.) and hence is not done.

..............................................
Question 10: What is the best confirmation point
            for interpolation?

For the usual continuous factor case, we note in
passing that the best (most efficient and highest
leverage) additional model-validation point that
may be added to a 2**k or 2**(k-p) design is at
the center point.  This center point augmentation
"costs" the experimentalist only one additional
run.  For example, for the k = 2 factor
(Temperature (300 to 350), and time (20 to 30))
experiment as discussed before, the usual 4-run
2***2 full factorial design:

        X1     X2      Y
         -      -      2
         +      -      4
         -      +      6
         +      +      8

may be augmented by the following 5-run 2**2 full
factorial design with a center point.  :

        X1     X2      Y
         -      -      2
         +      -      4
         -      +      6
         +      +      8
         0      0

Since - stands for -1 and + stands for +1, then it
is natural to code the center point as (0,0).
Using the recommended model, it is clear that the
predicted value Yhat at the center point will be

   Yhat(X1 = 0, X2 = 0) =
   Yhat = 5 + (1/2)(4*X2 + 2*X1 + 0*X1*X2)
   Yhat = 5 + (1/2)(4*(0) + 2*(0) + 0*(0)(0))
   Yhat = 5 + (1/2)(0)
   Yhat = 5 = Ybar

The design, data, and predicted values are thus:

        X1     X2      Y     Yhat
         -      -      2       2
         +      -      4       4
         -      +      6       6
         +      +      8       8
         0      0      ?       5  (confirmatory run)

The importance of the confirmatory run cannot be
overstated.  If the confirmatory run at the center
point yields a data value of, say, y = 5.1, then
(since the predicted value at the center is Yhat =
5 and since we know the model is perfect at the
corner points) that would give the analyst
considerable confidence in the quality of the
fitted model over the entire interior
(interpolation) domain.  On the other hand, if the
confirmatory run yielded a center point data value
quite different (e.g., Y = 7.5) from the center
point predicted value of Yhat = 5, then that would
prompt the analyst to NOT trust the fitted
model--even for interpolation purposes.  Hence
when our factors are continuous, a single
confirmatory run at the center point helps
immensely in assessing the range-of-trust for our
model.

<Replicated center points>
In practice, this center point value is frequently
repeated (2 reps, or even 3 reps) so as to not
only provide a reference point for assessing the
interpolatory power of the model at the center,
but also, the reps allow us to compute model-free
estimates of the natural error, which in turn
allows us a more rigorous method for computing the
uncertainty for individual coefficients in the
model, and for rigorously carrying out a
lack-of-fit test for assessing general model
adequacy.

..............................................
Question 11: How do we use the model for interpolation?

As for the mechanics of interpolation/
extrapolation itself, consider a continuation of
the prior k = 2 factor experiment.  Suppose
temperature T ranges from 300 to 350 and time t
ranges from 20 to 30, and the analyst can afford n
= 4 runs.  A 2**2 full factorial design is run.
Forming the coded temperature as X1 and the coded
time as X2, we have the usual:

        Temp    Time    X1     X2     Y
         300     30     -      -      2
         350     20     +      -      4
         300     30     -      +      6
         350     30     +      +      8

Graphically the design and data are as follows:

    30 +1 --------------
          |6          8|
     T    |            |
     i  X |            |
     m  2 |            |
     e    |            |
          |            |
          |2          4|
    20 -1 --------------
          -1    X1    +1
          300  Temp  350

As before, from the data, the average
"perfect-fit" prediction equation is

   Yhat = 5 + (1/2)(4*X2 + 2*X1 + 0*X1*X2)

We now pose the following typical interpolation
question:

   From the model, what is the predicted response
   at, say, temperature = 310 and time = 26?

In short:

   Yhat(T = 310, t = 26) = ?

To solve this problem, we first view the k = 2
design and data graphically, and note (via a ?) as
to where the desired (T = 310, t = 26)
interpolation point is:

    30 +1 --------------
          |6          8|
     T    |            |
     i  X |  ?         |
     m  2 |            |
     e    |            |
          |            |
          |2          4|
    20 -1 --------------
          -1    X1    +1
          300  Temp  350

The important next step is to convert the raw (in
units of the original factors T and t)
interpolation points into coded (in units of X1
and X2) interpolation points.  From the graph or
otherwise, we note that a linear translation
between T and X1, and between t and X2 yields

   T = 300 => X1 = -1
   T = 350 => X1 = +1

thus

   X1 = 0 is at T = 325


        |-------------|-------------|
       -1     ?       0            +1
       300   310     325           350

which in turn implies that

   T = 310 => X1 = -0.6

Similarly,

   t = 20 => X2 = -1
   t = 30 => X2 = +1

therefore,

   X2 = 0 is at t = 25

        |-------------|-------------|
       -1             0   ?        +1
       20             25 26        30

thus

   t = 26 => X2 = +0.2

Hence

     Yhat(T = 320, t = 26)
   = Yhat(X1 = -0.6, X2 = +0.2)
   = 5 + (1/2)(4*X2 + 2*X1 + 0*X1*X2)
   = 5 + (1/2)(4*(+0.2) + 2*(-0.6) + 0*(-0.6)(+0.2))
   = 5 + (1/2)(+0.8 - 1.2)
   = 5 + (1/2)(-0.4)
   = 5 + (-0.2)
   = 4.8

Thus

    30 +1 --------------
          |6          8|
     T    |            |
     i  X | 4.8        |
     m  2 |            |
     e    |            |
          |            |
          |2          4|
    20 -1 --------------
          -1    X1    +1
          300  Temp  350

..............................................
Question 12: How do we use the model for extrapolation?

For extrapolation, we do similarly.  For example,
the predicted value (denoted by ?) at (say) temperature T = 375,
and time t = 28 is indicated by the "?":

    30 +1 --------------
          |6          8|
     T    |            |     ?
     i  X |            |
     m  2 |            |
     e    |            |
          |            |
          |2          4|
    20 -1 --------------
          -1    X1    +1
          300  Temp  350

and is computed by

     Yhat(T = 375, t = 28)
   = Yhat(X1 = +2.0, X2 = +0.8)
   = 5 + (1/2)(4*X2 + 2*X1 + 0*X1*X2)
   = 5 + (1/2)(4*(+0.8) + 2*(+2.0) + 0*(+2.0)(+0.8))
   = 5 + (1/2)((3.2 + 4.0)
   = 5 + (1/2)(7.2)
   = 5 + 3.6
   = 8.6

and thus we have

    30 +1 --------------
          |6          8|
     T    |            |   8.6
     i  X |            |
     m  2 |            |
     e    |            |
          |            |
          |2          4|
    20 -1 --------------
          -1    X1    +1
          300  Temp  350

<Pseudo-data>
The predicted value from the modeling effort may
be viewed as pseudo-data-- an infinite amount of
"free" data without the experimental effort.  Such
"free" data can add tremendously to the insight
via the application of graphical techniques (in
particular, the contour plots of step 10), and can
add significant insight and understanding as to
the nature of the response surface relating Y to
the X's.

But again a final word of caution: the "pseudo
data" that results from the modeling process is
exactly that--pseudo-data-- it is NOT real data,
and so the model and the model's predicted values
must be validated by additional confirmatory
(real) data points.  As John Tukey remarked:

   "Models are to be trusted as real [that is
   generate predicted values and contour curves],
   but must always be verified" [that is, by
   additional confirmatory data points].

The rule of thumb is thus to take advantage of the
available and recommended model-building mechanics
for these 2-level designs, but do treat the
resulting derived model with an equal does of both
optimism and caution.

In summary, the motivation for model-building is
that it gives us insight into the nature of the
response surface along with the ability to do
interpolation and extrapolation.  Further, the
motivation for the use of the cumulative residual
standard deviation plot is that is serves as an
easy-to-interpret tool for determining such a good
and parsimonious model.

5. The Plot for the Defective Springs Data

Applying the |Effects| Plot to the defective
springs yields the following:

   <show plot>

6. How to Interpret

As discussed in detail under question 4 in the
Motivation section, the cumulative residual
standard deviation "curve" will characteristically
decrease left to right as we add more terms to the
model.  The incremental improvement (decrease)
tends to be large at the beginning (when important
factors are being added), but then the decrease
tends to be marginal at the end (as unimportant
factors are being added).

Including all terms would yield a perfect fit
(residual standard deviation = 0) but would also
result in an unwieldy model.  Including only the
first term (the average) would yield a simple
model (only one term!) but typically will fit
poorly.  Although a formal quantitative stopping
rule can be developed based on statistical theory,
a less-rigorous (but good) alternative stopping rule
that is graphical, easy to use, and highly
effective in practice is as follows:

   Keep adding terms to the model until the curve's
   "elbow" is encountered.  The "elbow point" is that
   value in which there is a consistent, noticeably
   shallower slope (decrease) in the curve.  Include
   all terms up to (and including) the elbow point
   (after all, each of these included terms decreased
   the residual standard deviation by a large
   amount).  Exclude any terms AFTER the elbow point
   (since all such successive terms decreased the
   residual standard deviation so slowly that the
   terms were "not worth the complication of
   keeping").

From the residual standard deviation plot for the
defective springs data, we note the following:

   1. The residual standard deviation (rsd) for
      the "baseline" model Yhat = Ybar = 71.25 is
      sres = 13.7.

   2. As we add the next term (X1), the rsd drops
      nearly 7 units (from 13.7 to 6.6).

   3. As we add yet another term (X1*X3), the rsd
      also drops another 3 units (from 6.6 to
      3.4).

   4. As we add another term (X2), the rsd goes
      from 3.4 to 1.5-- dropping about another 2
      units.

   5. When the next term is added (X3), the
      reduction in the rsd (from about 1.5 to 1.3)
      is negligible.

   6. Thereafter to the end, the total reduction
      in the rsd is from only 1.3 to 0.

In this case, the "kink" in the residual standard
deviation curve is at the X2 term.  Prior to that,
all added terms (including X2) reduced the rsd by
a large amount (7, then 3, then 2).  After the
addition of X2, the reduction to the rsd was small
(all less than 1): .2, then .8, then .5, then 0.

The final recommended model in this case thus
involves p = 4 items:

   1. the average
   2. factor X1
   3. the X1*X3 interaction
   4. factor X2

The fitted model thus takes on the form

   Yhat = average + (1/2)*(B1*X1 + B13*X1*X3 + B2*X2)

The least squares estimates for the coefficients
in this model are

   average = 71.25
   B1 = 23
   B13 = 10
   B2 = -5

The B1 = 23, B13 = 10, and B2 = -5 least squares
values are, of course, identical to the estimated
effects E1 = 23, E13 = 10, and E2 = -5 (= Yhat(+)
- Yhat(-1)) values as previously derived in step 7
of this recommended 10-step DEX analysis
procedure.

The final fitted model is thus

   Yhat = 71.25 + (1/2) * (23*X1 + 10*X1*X3 - 5*X2)

Applying this prediction equation to the 8 design
points yields: predicted values Yhat which are
close to the data Y, and residuals (Res = Y-Yhat)
which are close to zero:

        X1     X2     X3       Y    Yhat     Res
         -      -      -      67    67.25   -0.25
         +      -      -      79    80.25   -1.25
         -      +      -      61    62.25   -1.25
         +      +      -      75    75.25   -0.25
         -      -      +      59    57.25   +1.75
         +      -      +      90    90.25   -0.25
         -      +      +      52    52.25   -0.25
         +      +      +      87    85.25   +1.75

Computing the residual standard deviation (from
the formula given under question 2 in the
Motivation section):

   sres = sqrt(sum of squared residuals / (n-p))

where n = number of data points = 8, and p =
number of estimated coefficients (including the
average) = 4 yields

   sres = 1.54
   sres = 1.5 (if rounded to 1 decimal place).

This detailed res = 1.54 calculation brings us full
circle for 1.54 is the value given above
the X3 term on the cumulative residual standard
deviation plot.

7. Conclusions for the Defective Springs Data

The application of the Cumulative Residual
Standard Deviation Plot to the defective springs
data set results in the following conclusions:

   1. Good-fitting Parsimonious (constant + 3
      terms) Model:

         Y = 71.25 + (1/2) (23*X1 + 10*X1*X2 - 5*X2)

   2. Residual Standard Deviation for this Model
      (as a measure of the goodness-of-fit for
      this model):

          sres = 1.54

(control L) 
